% Source: sources/returns/S6_04_MAPS_MANIFOLDS/sections/manifolds.tex:8836–8955
\paragraph{Theorem D.3.1 (Proved: the global moving-zero family).}
Let $S\to\mathbb P^1$ be Engel's rational elliptic surface, with
fibres $IV^*,III,I_1$ at $0,1,\infty$, respectively. Put
$L=\mathcal O_S(P-O)$, $T=t_{6P}$ and
$J=L\otimes T^*L^{-1}\simeq\pi^*\mathcal O_{\mathbb P^1}(1)$,
with the direction of $T$ fixed in K13. We write $X^{\log}(u)$ for
the logarithmically transformed sphere, keeping it distinct from K13's
untwisted $X(u)$. Choose sections
$\phi_0,\phi_1$ of $\mathcal O(1)$ with divisors $\infty,0$,
respectively, normalized so that near infinity
$\phi_0=q\phi_1$, $q=1/t$. Set
$\phi_s=\pi^*(\phi_0-s\phi_1)$.
For every sufficiently small fixed $u\ne0$ there is a proper flat
holomorphic family $\mathcal Y_u\to\Delta_s$, with a holomorphic
map to $\mathbb P^1$, having the following properties. We perform
the same logarithmic transformations of multiplicities $3,4$ and
the same admissible sphere twist data as for the chosen $X^{\log}(u)$.
The central threefold $Y_{u,0}$ has exactly one ordinary double point
and is the contraction of the diagonal $(-1,-1)$-curve of $X^{\log}(u)$.
Every fibre $Y_{u,s}$ for $s\ne0$ is smooth. Its local change at the
contracted point is
\[
 xy=u(zv-s),\qquad q=zv.
\]
This is the conifold transition described in Engel's Remark 8.5.
It is transverse to, and different from, K13's $u\to0$ degeneration.
\begin{proof}
The equality defining $J$ and its normalization are checked in
K13, A.1.2, against Engel's Lemma 2.1 and Corollary 2.2. We retain
that equality, but allow the section of $J$ to vary. Here are the
charts and their identifications. Choose Hermitian norms and
$0<r<1/4$. On the compact surface $S$, after restricting both
parameters there is a uniform inequality
$|u|\|\phi_s(x)\|<r^2$. Take a core
\[
 C=\{(x,a,s):a\in L_x,\ r<\|a\|<r^{-1}\}
\]
and a node chart
\[
 N=\{(x,a,b,s):a\in L_x,\ b\in L^{-1}_{Tx},\
           a\otimes b=u\phi_s(x),\ \|a\|,\|b\|<2r\}.
\]
On the collar $r<\|a\|<2r$, identify $N$ with $C$ by
$(x,a,b,s)\mapsto(x,a,s)$. On $r<\|b\|<2r$, identify it by
$(x,a,b,s)\mapsto(Tx,b^{-1},s)$. These maps and their inverses are
holomorphic; on the first collar $b=u\phi_s/a$, and on the second
$a=u\phi_s b^{-1}$. The two node collars are disjoint by the
uniform inequality. A change of local line frames multiplies both
sides of the node equation by the same unit, so the construction is
intrinsic.

The quotient is Hausdorff. At each end of an attaching collar one
chart contains the limiting circle and the other omits it; there is
no pair of distinct unpaired limits of the gluing graph. The two
node collars are separated, so no new chain of identifications
changes this conclusion. For properness, the closed core
$3r/2\leq\|a\|\leq2/(3r)$ and the closed node subset
$\|a\|,\|b\|\leq7r/4$ cover the quotient. Over any compact set
of parameters their preimages are compact because $S$ is compact.
This proves properness, also for the joint map to
$\mathbb P^1\times\Delta_s$.

On the core, flatness over $s$ is immediate. On the node chart, the
local rings are hypersurface rings over the smooth surface $S$;
the equation $ab=u\phi_s$ is not divisible by $s-s_0$, so
multiplication by $s-s_0$ is injective. A coherent module over the
one-dimensional regular local parameter ring is flat precisely
when it is torsion free; hence the family is flat. Alternatively,
for the joint map to $(t,s)$ the fibres all have pure dimension
two, and the Cohen--Macaulay regular-sequence criterion gives
flatness over the smooth two-dimensional base.

Away from the zero of $\phi_s$, the collar relation gives exactly
$L^\times/\langle u\phi_s\rangle$, with the contracting generator
covering $T^{-1}$. The uniform norm inequality ensures that every
orbit meets a core or node representative, so this is an actual
biholomorphic comparison. Near $0$ and $1$ the ratio
$\phi_s/\phi_0$ is a holomorphic unit, uniformly for small $s$.
A local logarithm changes the scalar period by a holomorphic
function. Thus the same good-reduction torus cover and the same
finite affine translation data extend holomorphically in $s$.
Their finite cyclic actions stay free, by the unchanged primitive
invariant translation coordinate. We replace the two corresponding
pieces by these smooth logarithmic-transform models and glue on
the punctured overlaps by the just-described period comparison.
This is precisely the original finite gluing when $s=0$.
Properness and flatness persist, since the replacements are proper
families with the same compact annular boundaries and are locally
finite free quotients of smooth families.

At the node of $S_\infty$ write $q=zv$. A unit change in a line
coordinate gives $xy=u(zv-s)$. The relative gradient in
$(x,y,z,v)$ is $(y,x,-uv,-uz)$, so the only possible critical point
is the origin, and it lies on the fibre exactly when $s=0$.
At a smooth point of a zero fibre of $\phi_s$, its zero is simple,
so $xy=u(q-s)$ is smooth. These observations exhaust the singular
points: the other nodes are absent because the elliptic surface
has only its stated three exceptional fibres, and the two finite
ones were replaced by the smooth free quotients. At $s=0$ the
quadratic form $xy-uzv$ is nondegenerate and gives one ordinary
double point. Its small resolution is the diagonal subdivision of
K13. In one chart write $x=az$, $uv=ay$; in the other write
$z=bx$, $y=buv$, with $b=a^{-1}$. The two normal coordinates
transform by multiplication by $a$, giving
$\mathcal O(-1)\oplus\mathcal O(-1)$ over the exceptional
$\mathbb P^1$. There is one such diagonal modulo the deck periods.
The comparison with K13 identifies its small resolution with the
chosen $X^{\log}(u)$.

Before the finite logarithmic replacements, both constructions sit in
the global core/node plumbing with the explicit two-parameter equation
$xy=u(zv-s)$. The checkpoint follows the slice $s=0$ and varies
$u$ towards zero; the present construction fixes $u\ne0$ and varies
$s$. At $u=0$ the node equation is $xy=0$ and recovers the same
untwisted pinched projective bundle for every $s$. The finite
logarithmic replacements used for the sphere have only been constructed
here for $u\ne0$; no extension of those replacements across $u=0$
is asserted. The common cusp chart and the untwisted global plumbing
are the precise comparison with K13.
\end{proof}
